\documentclass{article}

\usepackage{iclr2027_conference,times}
\input{math_commands.tex}

\usepackage{amsmath}
\usepackage{amssymb}
\usepackage{booktabs}
\usepackage{graphicx}
\usepackage{multirow}
\usepackage{url}
\usepackage{hyperref}
\usepackage{xcolor}

% Primary contribution: World-State Chunk Retention (future-use default).
\newcommand{\method}{World-State CR}
\newcommand{\eg}{e.g.,\ }
\newcommand{\ie}{i.e.,\ }

\title{World-State Chunk Retention for Long-Horizon\\
Video World Models}

\author{Anonymous Authors}

% Keep \iclrfinalcopy commented for anonymous submission.
%\iclrfinalcopy

\begin{document}

\maketitle

\begin{abstract}
Autoregressive video world models store long-horizon evidence in key--value
(KV) caches.  Full history preserves revisits but breaks real-time budgets;
a sliding window is fast but forgets previously seen viewpoints.
We argue that the dominant prior practice for deployable world-model
inference is sliding-window KV, and propose \textbf{Chunk Retention} (CR):
a three-tier post-generation policy that always keeps a sink and a recent
window, while ranking older chunks into a capacity-limited archive.
Our default scorer, \textbf{World-State CR}, is an 11-D pose-aware
ChunkSelector trained against a \emph{Future Coverage Oracle}: a discounted
mix of future attention mass and pose/frustum reuse, with corrected
world-state feature formulas (\texttt{world\_state.v2}).
Heuristic CR, a 5-D Learned CR, and an attention-mass World-State ablation
share the same tier budget.  On RealCam-Vid we evaluate revisit fidelity
with the WorldKV protocol (PSNR / SSIM / LPIPS / FID) together with
throughput and peak memory, comparing against Sliding Window and
training-free WorldKV retrieval.  Token-level saliency pruning is optional
and \emph{not} required for the main claim.
\end{abstract}

\section{Introduction}
\label{sec:intro}

Interactive video world models generate camera- or action-conditioned
streams autoregressively~\citep{huang2025selfforcing,he2025matrix,
robbyant2026lingbot}.  Causal attention KV caches provide continuity, but
context grows with rollout length.  In practice, the widely usable baseline
for long rollouts is a \textbf{sliding window}: bound active KV, discard
evicted chunks.  This is cheap, yet fails when the camera returns to a past
viewpoint---the original scene evidence is gone.

Recent work explores training-free KV retrieval and
compression~\citep{yi2026worldkv}, gated long-video
caches~\citep{meng2026tether}, and related temporal compression
policies~\citep{luo2026future,samuel2026fastar}.  We ask a narrower systems
question for \emph{chunked} LingBot-World-style generators:
\begin{quote}
\emph{Under a fixed active KV budget, which historical chunks should remain
on device after each new chunk is generated?}
\end{quote}

Our answer is \textbf{tiered Chunk Retention} with a
\textbf{world-state} ranking signal trained for \emph{future coverage}
(Figure~\ref{fig:overview}).  Unlike host-side demote/retrieve graphs or
heavy token cascades, CR is a post-generation retention policy: generate a
chunk, then decide what stays in the active attention context.  Dropped
archive chunks are not recovered later in the rollout---budget discipline
is explicit.

\begin{figure}[t]
  \centering
  \fbox{\begin{minipage}[c][3.2cm][c]{0.94\linewidth}
  \centering
  \textbf{Tiered Chunk Retention}\\[2mm]
  Sink (always) \;/\; Recent (FIFO) \;/\; Archive (ranked).\\[1mm]
  Archive utility: World-State CR (11-D, Future Coverage Oracle).
  \end{minipage}}
  \caption{Prior deployable baseline keeps only a sliding window.
  \method{} keeps sink\,+\,recent\,+\,a state-ranked archive under the same
  token budget.  The default ranker is trained for future pose/attention
  coverage rather than instantaneous attention mass.}
  \label{fig:overview}
\end{figure}

\paragraph{Contributions.}
(1)~A simple three-tier CR skeleton (sink / recent / archive) as the
retention interface for chunked world models.
(2)~\textbf{World-State CR} (default): an 11-D ChunkSelector
(\texttt{world\_state.v2}) trained with a Future Coverage Oracle
(\texttt{future\_use\_v1}), checkpoint \texttt{selector\_ws\_future\_v1.pt}.
(3)~Controlled ablations (Heuristic CR, Learned CR, attention-mass
World-State v1) and comparisons to Sliding Window and WorldKV under a
WorldKV-style revisit protocol on RealCam-Vid.

\section{Method}
\label{sec:method}

\subsection{Tiered Chunk Retention}
LingBot-World-Fast generates video in latent chunks of size $C{=}3$.  After
chunk $t$ is written into the KV cache, history is partitioned into:
\begin{itemize}
  \item \textbf{Sink}: the first $S$ chunks (default $S{=}1$), never evicted;
  \item \textbf{Recent}: the last $R$ chunks (default $R{=}1$), FIFO;
  \item \textbf{Archive}: older candidates ranked by utility; keep
        $\max\!\big(K_{\min},\,\lceil \rho N\rceil\big)$ under the attention
        token budget.
\end{itemize}
Default tier ratio is $S{:}R{:}K_{\min}=1{:}1{:}2$ with keep ratio
$\rho{=}0.5$ (selected by a RealCam-Vid ratio sweep).  Ranking happens
\emph{after} generation; evicted archive chunks are discarded for the
remainder of the rollout.

\subsection{Archive scorers}
\paragraph{Heuristic CR (ablation).}
Utility is a peak-mean camera/action motion score, fused with latent
residual rescue for static-camera dynamic content.  No learned weights.

\paragraph{Learned CR (ablation).}
A tiny MLP over a 5-D feature
$(\mathrm{motion},\,\mathrm{age},\,\mathrm{cam\text{-}angular},\,
\mathrm{key\text{-}centroid},\,\mathrm{value\text{-}norm})$ predicts
utility.  Checkpoint: \texttt{selector\_all4.pt} (schema
\texttt{learned.v0}).

\paragraph{World-State CR (default).}
\label{sec:ws-cr}
We use an 11-D geometric feature layout (schema \texttt{world\_state.v2}):
motion norm; SE(3) translation / rotation geodesic to the current pose;
camera-forward overlap; sparse frustum-overlap proxy (pose\,+\,intrinsics,
no depth); relative motion and reachability over elapsed chunks;
key-centroid affinity; value norm; soft revisit count; time since last
observed.  Relative to the frozen v1 formulas, v2 corrects the
reachability, relative-motion, and time-since-last-observed normalizers
so that long-horizon rollouts do not saturate.

A small MLP (hidden 32, dropout 2) maps standardized features to a scalar
utility.  Checkpoint: \texttt{selector\_ws\_future\_v1.pt}.

\paragraph{Future Coverage Oracle.}
Rather than ranking chunks by instantaneous attention mass at time $t$, we
train against a \emph{future survival} target over horizon $H$ (default
$H{=}8$).  For candidate chunk $i$ at step $t$,
\begin{equation}
  y_{t,i}
  =
  \sum_{h=1}^{H}
  \gamma^{h-1}
  \Big[
    \alpha\,\mathrm{AttnMass}_{t+h}(i)
    +
    (1-\alpha)\,\mathrm{PoseReuse}_{t+h}(i)
  \Big],
\end{equation}
with defaults $\gamma{=}0.9$ and $\alpha{=}0.5$.
$\mathrm{PoseReuse}$ is a cheap frustum\,+\,forward-overlap proxy between
the stored chunk and a future camera pose.  Labels are converted offline
from attention-mass oracle dumps (\texttt{train\_selector.py
--label\_type future\_use\_v1 --feature\_schema v2}); inference remains
fully causal and uses only the 11-D features above.

\paragraph{World-State CR v1 (ablation).}
The earlier attention-mass checkpoint \texttt{selector\_ws\_v1.pt}
(\texttt{world\_state.v1}) shares the same CR skeleton and 11-D layout but
uses the frozen feature formulas and instantaneous attention-mass labels.
We report it as an ablation against the default future-use ranker
(\texttt{world\_state\_cr\_v1} in code).

\subsection{Optional token pruning (not the main claim)}
Saliency-weighted token pruning inside archive chunks (SWTP) can further
compress tokens after promotion.  Empirically, residual saliency is a weak
proxy for revisit utility, so we \textbf{do not} treat SWTP---or the
CR\,+\,SWTP product name ``MoSaiC''---as central to the paper narrative.
World-State CR stands alone as the primary method.

\section{Related Work}
\label{sec:related}

World models with persistent
memory~\citep{robbyant2026lingbot,xiao2025worldmem,he2025matrix} motivate
revisit consistency.  WorldKV~\citep{yi2026worldkv} stores evicted KV in a
bank and retrieves by camera/action similarity, optionally compressing
tokens at store time.  Our CR instead \emph{irreversibly} sparsifies the
active context with a learned world-state ranker under a fixed on-device
budget---complementary to banked retrieval.  Long-video KV
managers~\citep{meng2026tether,luo2026future,samuel2026fastar} target drift
or temporal redundancy; we target viewpoint revisits in camera-conditioned
world rollouts.

\section{Experiments}
\label{sec:exp}

\paragraph{Setup.}
Backbone: LingBot-World-Fast (camera-conditioned).  Benchmark: RealCam-Vid
default battery (\texttt{default\_all}: 24 loop\,+\,40 random clips;
primary memory results on \texttt{default\_loop}).  Matched seeds and
resolution; CR uses $\texttt{max\_attention\_size}{=}47000$ and tier ratio
$1{:}1{:}2$.  Baselines: Sliding Window (\texttt{window}) and official
WorldKV (\texttt{use\_retrieval}, $\texttt{retrieval\_frames}{=}18$,
store-time compression, $N{=}2{\times}U{=}4$ on 8$\times$H20).

\paragraph{Metrics (WorldKV protocol).}
For each GT pose trajectory we detect revisit frames (SE(3) radius $0.15$,
minimum gap $30$ frames, at most $64$ pairs per clip) and compare each
revisit frame to its first-visit frame with \textbf{PSNR / SSIM / LPIPS};
\textbf{FID} between revisit and first-visit sets; plus wall-clock
\textbf{FPS}, mean attention \textbf{context} tokens, and \textbf{peak}
GPU memory.
Heuristic / Learned / World-State v1 are reported as ablations under the
same CR budget.

\paragraph{Main results.}
Table~\ref{tab:main-loop} reports the WorldKV revisit protocol on
\texttt{default\_loop} (24 clips; all have revisit pairs).
The default \method{} is the future-use ranker
(\texttt{selector\_ws\_future\_v1.pt} / \texttt{world\_state.v2});
the attention-mass World-State checkpoint
(\texttt{selector\_ws\_v1.pt} / \texttt{world\_state.v1}) is shown as an
ablation under the same CR tiers.
Relative to Sliding Window, default \method{} improves revisit PSNR /
SSIM / LPIPS / FID while cutting context tokens by $\sim$48\% and peak
memory by $\sim$28\%, at higher FPS.  Future-use training further reduces
FID versus the attention-mass World-State ablation (61.4 vs.\ 66.2) at
matched efficiency.

\begin{table}[t]
  \centering
  \caption{WorldKV revisit protocol on RealCam-Vid \texttt{default\_loop}
  (24 clips).  Default \method{} = future-use / \texttt{world\_state.v2};
  World-State v1 = attention-mass ablation.}
  \label{tab:main-loop}
  \begin{tabular}{lccccccc}
  \toprule
  Method & PSNR$\uparrow$ & SSIM$\uparrow$ & LPIPS$\downarrow$ & FID$\downarrow$
         & FPS$\uparrow$ & ctx & peakGB \\
  \midrule
  Window & 9.830 & 0.300 & 0.726 & 67.0 & 1.18 & 41471 & 54.7 \\
  World-State v1 & 10.026 & 0.327 & 0.719 & 66.2
         & \textbf{1.47} & \textbf{21489} & \textbf{39.2} \\
  \method{} (default) & \textbf{10.138} & \textbf{0.328} & \textbf{0.713}
         & \textbf{61.4} & \textbf{1.47} & \textbf{21489} & \textbf{39.2} \\
  \bottomrule
  \end{tabular}
\end{table}

\paragraph{Efficiency takeaway.}
Both World-State variants retain $\sim$4 chunks under the $1{:}1{:}2$ tier
policy and share the same FPS / ctx / peakGB; gains of the default
future-use ranker are therefore attributable to the Future Coverage Oracle
and \texttt{world\_state.v2} formulas rather than a different budget.
Sliding Window is neither the most faithful nor the most efficient once a
world-state archive is available.

\paragraph{Additional comparisons.}
Official WorldKV and Heuristic / Learned CR ablations use the same
protocol (\texttt{bench/eval\_worldkv\_memory.py}); CLI names are listed
in Section~\ref{sec:code}.

\section{Code Alignment}
\label{sec:code}

\begin{center}
\begin{tabular}{ll}
\toprule
CLI / method & Role \\
\midrule
\texttt{window} & Sliding-window baseline \\
\texttt{heuristic\_cr} & Heuristic CR (ablation) \\
\texttt{learned\_cr} & Learned CR, 5-D (ablation) \\
\texttt{world\_state\_cr} & \method{} (default; future-use) \\
\texttt{world\_state\_cr\_v1} & Attention-mass World-State ablation \\
\texttt{worldkv} batch & Official WorldKV baseline \\
\bottomrule
\end{tabular}
\end{center}

Default checkpoint: \texttt{assets/selectors/selector\_ws\_future\_v1.pt}.
The former experimental name \texttt{world\_state\_cr\_future} is retained
as a back-compat alias of \texttt{world\_state\_cr}.

\section{Conclusion}
\label{sec:conclusion}

We reframe sparse KV for video world models around
\textbf{tiered Chunk Retention} and a default \textbf{World-State} archive
ranker trained for \textbf{future coverage}.  Sliding Window remains the
practical prior; Heuristic, Learned, and attention-mass World-State are
ablations; WorldKV is the natural retrieval baseline.  Optional token
pruning is secondary.  The release matches this terminology
(\texttt{world\_state\_cr} $\equiv$ future-use default) and the
WorldKV-style revisit evaluation protocol.

\bibliography{iclr2027_conference}
\bibliographystyle{iclr2027_conference}

\end{document}
